% Part: lambda-calculus
% Chapter: introduction
% Section: reduction

\documentclass[../../../include/open-logic-section]{subfiles}

\begin{document}

\olfileid{lam}{int}{red}
\olsection{లాంబ్డా పదాల తగ్గింపు}

లాంబ్డా పదాలతో ఏమి చేయవచ్చు? వాటిని సరళీకరించవచ్చు. $M$, $N$
లాంబ్డా పదాలూ, $x$ ఒక చరమూ అనుకుందాం. ప్రతిస్థాపన వల్ల~$N$లోని
స్వేచ్ఛా చరాలు బద్ధమయ్యే ప్రమాదం ఉంటే, అందుకు కారణమయ్యే $M$లోని
బద్ధ చరాల పేర్లను ముందే మారుస్తాం. ఆ తరువాత $M$లో~$x$ స్థానంలో $N$ను
ప్రతిస్థాపించి వచ్చే పదాన్ని $\Subst{M}{N}{x}$తో సూచిస్తాం. ఉదాహరణకు,
\[
\Subst{(\lambd[w][xxw])}{yyz}{x} = \lambd[w][(yyz)(yyz)w].
\]

\begin{digress}
ప్రతిస్థాపనకు ఇతర సంకేతనాలు $[N/x]M$, $[x/N]M$, అలాగే $M[x/N]$.
జాగ్రత్త!
\end{digress}

సహజంగా చూస్తే, $(\lambd[x][M])N$, $\Subst{M}{N}{x}$లకు ఒకే
అర్థం ఉంటుంది. మొదటి పదం స్థానంలో రెండో పదాన్ని పెట్టడాన్ని
\emph{$\beta$-సంకోచనం} అంటారు. $(\lambd[x][M])N$ను
\emph{రెడెక్స్} అనీ, $\Subst{M}{N}{x}$ను దాని
\emph{సంకోచిత ఫలితం} (కాంట్రాక్టమ్) అనీ అంటారు. సాధారణంగా, పదం
$P$లోని ఒక ఉపపదానికి $\beta$-సంకోచనం చేసి మొత్తం పదాన్ని~$P'$గా మార్చగలిగితే,
$P$ \emph{ఒక దశలో $P'$కు $\beta$-తగ్గుతుంది} అంటాం. దీన్ని
$P \redone P'$ అని రాస్తాం. $P$ నుంచి ఒకదశ తగ్గింపులను వరుసగా చేసి
(లేదా ఏ దశా లేకుండానే)~$P'$ను పొందగలిగితే, $P$ \emph{$P'$కు
$\beta$-తగ్గుతుంది} అంటాం. దీని సంకేతం $P \red P'$. ఇకపై $\beta$-తగ్గింపు
సాధ్యం కాని పదాన్ని \emph{$\beta$-తగ్గించలేనిది} లేదా \emph{$\beta$-నియత రూపం}
అంటాం. సందర్భం స్పష్టంగా ఉన్నప్పుడు ``$\beta$-తగ్గుతుంది'' మొదలైన వాటికి
బదులు ``తగ్గుతుంది'' అని మాత్రమే అంటాం.

కొన్ని ఉదాహరణలను పరిశీలిద్దాం.
\begin{enumerate}
\item మనకు
\begin{align*}
(\lambd[x][xxy]) \lambd[z][z] & \redone (\lambd[z][z])(\lambd[z][z]) y \\
& \redone (\lambd[z][z]) y \\
& \redone y.
\end{align*}
\item ఒక పదాన్ని ``సరళీకరించడం'' దాన్ని మరింత సంక్లిష్టంగా చేయవచ్చు:
\begin{align*}
(\lambd[x][xxy])(\lambd[x][xxy]) & \redone (\lambd[x][xxy])(\lambd[x][xxy])y \\
& \redone (\lambd[x][xxy])(\lambd[x][xxy])yy \\
& \redone \dots
\end{align*}
\item అది ఒక పదాన్ని మార్చకుండా కూడా ఉంచవచ్చు:
\[
(\lambd[x][xx])(\lambd[x][xx]) \redone (\lambd[x][xx])(\lambd[x][xx]).
\]
\item కొన్ని పదాలను ఒకటి కంటే ఎక్కువ విధాలుగా కూడా తగ్గించవచ్చు.
  ఉదాహరణకు,
\[
(\lambd[x][(\lambd[y][yx]) z)] v \redone (\lambd[y][yv]) z
\]
అనేది అన్నిటికంటే బయట ఉన్న ప్రయోగాన్ని సంకోచిస్తే వస్తుంది.
\[
(\lambd[x][(\lambd[y][yx]) z)] v \redone (\lambd[x][zx]) v
\]
అనేది అన్నిటికంటే లోపల ఉన్న ప్రయోగాన్ని సంకోచిస్తే వస్తుంది. ఈ ఉదాహరణలో
రెండు ఫలితాలూ ఇంకో తగ్గింపు తరువాత ఒకే పదం~$zv$ను ఇస్తాయి.
\end{enumerate}

చివరి ఉదాహరణలో రెండు దారులూ ఒకే ఫలితాన్ని ఇవ్వడం యాదృచ్ఛికం కాదు.
ఇది లాంబ్డా కలనశాస్త్రంలోని ``చర్చ్--రోసర్ లక్షణం'' అనే లోతైన, ముఖ్యమైన లక్షణాన్ని
చూపిస్తుంది.

\end{document}
